\chapter{Verification}\label{r03:ch:verification}

\emph{How each change was shown to compute the same results before its speed was measured.}

Every change was checked for correctness before its speed counted. Five kinds of check were
used, and \cref{r03:fig:vermap} shows which change each of them covered:
\begin{enumerate}
  \item the \ps{} test suite, on the laptop and natively on the board;
  \item waLBerla's scenario and halo-exchange tests, and \term{mutation checks}: the code under
        test is broken on purpose, and the test must then fail;
  \item the application's checksum, which must agree between variants and between different
        divisions of the domain into blocks;
  \item QEMU runs of the lab kernels on the laptop;
  \item the checksums printed by the board runs.
\end{enumerate}
Runs marked ``repeated for this report'' were made again on the committed code (\rcommitps,
\rcommitwl) while this report was written.

\begin{figure}[htbp]
  \centering
  \resizebox{\linewidth}{!}{% Which check covered which change (chapter 8). Filled dot: done, ring: pending, dash: not applicable.
\begin{tikzpicture}[every node/.style={font=\scriptsize},
    d/.style={circle,fill=r03landed!85!black,minimum size=2.2mm,inner sep=0pt},
    p/.style={circle,draw=r03pending!90!black,thick,minimum size=2.2mm,inner sep=0pt},
    n/.style={label={center:{\textcolor{psgray}{--}}}}]
  \def\rthreeColW{2.05}
  \foreach \t [count=\j] in {{\ps{} tests\\(laptop)},{RVV execution\\tests (board)},{QEMU: lab\\kernels},
      {x86: application,\\scenario tests},{mutation\\checks},{board: end to end,\\lab acceptance}} {
    \node[align=center,font=\scriptsize\bfseries,text width=19mm] at (2.9+\j*\rthreeColW,0.45) {\t};
  }
  \foreach \row/\name/\marks [count=\i] in {
      1/{software prefetch (\ref{r03:sec:ps-prefetch})}/{d,d,n,n,n,d},
      2/{AA streaming (\ref{r03:sec:wl-sweepgen}, \ref{r03:sec:wl-halo})}/{n,n,d,d,d,d},
      3/{.vf operands and FMA (\ref{r03:sec:ps-selector})}/{d,d,d,n,d,d},
      4/{row pointers, \texttt{in\_register} (\ref{r03:sec:ps-rowptr})}/{d,d,d,n,d,d},
      5/{halo one slot at a time (\ref{r03:sec:wl-halo})}/{n,n,n,d,d,p},
      6/{no density/velocity output (\ref{r03:sec:wl-sweepgen})}/{n,n,d,d,n,p},
      7/{split AA, generated (\ref{r03:sec:wl-split})}/{n,n,d,d,n,p}} {
    \pgfmathsetmacro{\y}{-\row*0.5}
    \node[anchor=west,text width=44mm] at (-0.25,\y) {\name};
    \foreach \m [count=\j] in \marks {
      \node[\m] at (2.9+\j*\rthreeColW,\y) {};
    }
  }
  \draw[psgray!50] (-0.25,-0.18) -- (15.4,-0.18);
  \node[anchor=west,r03label] at (-0.25,-4.05)
    {\tikz\fill[r03landed!85!black] (0,0) circle (1.1mm); done \quad
     \tikz\draw[r03pending!90!black,thick] (0,0) circle (1.1mm); pending (measurement queued) \quad
     -- not applicable};
\end{tikzpicture}
}
  \caption[Which check covered which change]{Which check covered which change. Numbers in brackets
    are the sections that describe the change.}\label{r03:fig:vermap}
\end{figure}

Changes to \ps{} reached the board through a guarded session script
(\cref{r03:fig:jobflow}). It installed the changed files, ran \ps{}'s RVV tests natively on the
X60, and restored the previous files if any test failed. Only then did it build and measure.

\begin{figure}[htbp]
  \centering
  % A guarded session script (misc/x60-lab/queue_e2e_cg.sh; queue_lab3.sh and queue_e2e_halo.sh follow
% the same pattern): changed files stay installed only if the board's tests pass. Right: the keep rule
% that queue_e2e_halo.sh applied to the row-pointer change of run C.
\begin{tikzpicture}[every node/.style={font=\scriptsize},node distance=3mm and 6mm,
    r03step/.append style={font=\scriptsize}]
  \node[r03info,r03step=42mm] (s0) {laptop: pack the changed files into an archive and copy it to the board};
  \node[r03step=42mm,below=of s0] (s1) {start in a detached \texttt{screen}; wait until the previous session script has finished};
  \node[r03step=42mm,below=of s1] (s2) {back up the board's copies of the files; unpack the archive};
  \node[r03decide,below=6mm of s2] (g) {\ps{} RVV tests\\on the board\\pass?};
  \node[r03no,r03step=42mm,right=12mm of g] (r) {restore the backup, remove new files, stop};
  \node[r03step=42mm,below=6mm of g] (s3) {delete the generated sources, so that they are regenerated};
  \node[r03step=42mm,below=of s3] (s4) {for f32 and f64: configure, build every variant};
  \node[r03step=42mm,below=of s4] (s5) {run each variant 3$\times$ on core 2, $64^3$ cells, 200 steps};
  \node[r03ok,r03step=42mm,below=of s5] (s6) {log: best MLUP/s and checksum of each variant};
  \foreach \u/\v in {s0/s1,s1/s2,s2/g,s3/s4,s4/s5,s5/s6} \draw[r03arrow] (\u) -- (\v);
  \draw[r03arrow] (g) -- node[above] {no} (r);
  \draw[r03arrow] (g) -- node[right] {yes} (s3);

  % the keep rule for the row-pointer change
  \node[r03decide,right=14mm of s4] (k) {AA + prefetch gained\\$\ge 2$\,\% in f32\\and in f64?};
  \node[r03box,text width=38mm,font=\scriptsize,above=7mm of k] (kk) {next session script, before it installs its own change};
  \node[r03box,text width=26mm,font=\scriptsize,below left=8mm and -6mm of k] (keep) {keep the \ps{} change};
  \node[r03box,text width=26mm,font=\scriptsize,below right=8mm and -6mm of k] (rev) {restore the previous \ps{} files};
  \draw[r03arrow] (kk) -- (k);
  \draw[r03arrow] (k) -- node[left] {yes} (keep);
  \draw[r03arrow] (k) -- node[right] {no} (rev);
\end{tikzpicture}

  \caption[A guarded session script]{A guarded session script: changed files stay installed
    only if the tests pass on the board. The right-hand part is the rule that decided whether
    the row-pointer change was kept (\cref{r03:sec:res-cg}).}\label{r03:fig:jobflow}
\end{figure}

\section{pystencils}\label{r03:sec:ver-ps}

\begin{table}[htbp]
  \centering\small
  \caption{\ps{} test runs.}\label{r03:tab:ver-ps}
  \begin{tabularx}{\linewidth}{@{}l>{\raggedright\arraybackslash}Xl@{}}
    \toprule
    \textbf{Where, when} & \textbf{Scope} & \textbf{Result}\\
    \midrule
    laptop, repeated for this report & \file{tests/nbackend} and \file{tests/codegen} at
      \texttt{a81dc547} & 435 passed, 15 skipped\\
    board, with software prefetch & RVV target: code-generation and execution tests, and
      \file{test_vectorization.py}; compiled and run natively on the X60 & 130 passed\\
    board, with the new instruction selection & the same (second attempt, see
      \cref{r03:sec:ver-bugs}) & 145 passed\\
    board, with row pointers and \code{in_register} & the same & 148 passed\\
    \bottomrule
  \end{tabularx}
\end{table}

Of the 15 tests skipped on the laptop, 13 need RVV hardware (those in
\file{test_rvv_execution.py}), one needs a GPU, and one needs the Python package \code{islpy}.
The board runs are in the logs \file{pytest-rvv.log}, \file{pytest-vf2.log} and
\file{pytest-cg.log} (Appendix~\ref{r03:app:evidence}).

\section{waLBerla}\label{r03:sec:ver-wl}

\paragraph{Scenario and halo-exchange tests.} In the laptop's x86 test build (GCC~16, with MPI):
\begin{shcode}
ctest -R "LbmScenarios|HaloExchange|CommunicationHiding"
\end{shcode}
\textbf{24 of 24 passed} at \texttt{efa2b70c} (repeated for this report, after regenerating the
scenario sources, which use \code{LbmBulk}). The tests are the LBM scenarios FullyPeriodic,
AAMatchesPull, MirroredHalfChannel and FreeSlipPipe, and the v8 core's tests
\code{testHaloExchangeInterfacesHost}, \code{testHostFieldPackInfo},
\code{testStreamPullPackInfo} and \code{testCommunicationHiding}, each with 1, 2 and 4
processes.

\paragraph{Mutation checks.} A test that cannot fail proves nothing. The AA halo exchange was
therefore broken on purpose, in three ways:
\begin{itemize}
  \item Without the AA push exchange's check whether a ghost slot was written by the sender,
        AAMatchesPull fails with 4 processes. It passes with 1 and 2, where the unpack order
        happens to let the correct value win: exactly the hazard the check removes.
  \item Sending the direct instead of the inverse slots in the AA pull exchange's \code{pack}
        fails with 2 and 4 processes.
  \item After the slot-by-slot rewrite: with \code{writtenGhosts} returning the whole ghost layer
        (no slots skipped), AAMatchesPull fails with 4 processes, the only case with edges
        between different blocks.
\end{itemize}

\section{The application's checksum on x86}\label{r03:sec:ver-app}

The application prints the sum of $|u|^2$ after the run (\cref{r03:sec:app}). In double
precision all variants must print exactly the same value, for any division of the same domain
into blocks: scalar and vector, pull and AA, single-pass and split, with and without output in
the timed steps.

\begin{table}[htbp]
  \centering\small
  \caption{Checksum after N time steps (2 warm-up steps), laptop x86 build at \texttt{efa2b70c},
    f64, vector variants with x86 AVX instructions (repeated for this report). Each value is
    printed, to all 12 digits, by all five variants \texttt{scalar\_cse}, \texttt{scalar\_cse\_aa},
    \texttt{explicit\_cse}, \texttt{explicit\_cse\_aa} and \texttt{explicit\_cse\_aa\_split}, each
    with \texttt{SWEEPGEN\_LBM\_OUTPUT=every} and \texttt{=last}: 60 runs.}\label{r03:tab:ver-app}
  \begin{tabular}{@{}lrr@{}}
    \toprule
    \textbf{Processes $\times$ block} & \textbf{steps} & \textbf{checksum}\\
    \midrule
    $1\times64^3$ and $8\times32^3$ & 20 & 395.489692731\\
    $1\times64^3$ and $8\times32^3$ & 21 & 395.575509355\\
    $2\times24^3$ & 101 & 40.7081176884\\
    $6\times24^3$ & 101 & 125.363104502\\
    \bottomrule
  \end{tabular}
\end{table}

The table covers both kinds of halo traffic, the periodic copy of one block onto itself
(\code{localCopy}) and the exchange between 2, 6 and 8 processes (\code{pack}/\code{unpack}),
after even and after odd numbers of steps. The same agreement was recorded when AA streaming was
first implemented, before the slot-by-slot rewrite and before the split kernel and the kernels
without output existed.

\section{QEMU runs of the lab kernels}\label{r03:sec:ver-qemu}

The lab driver with one kernel, cross-compiled statically with GCC~16 and run under
\code{qemu-riscv64 -cpu rv64,v=true,vlen=256,zicbop=true}, prints the three checksums of
\cref{r03:sec:lab}: $|u|^2$ sum, mass and weighted PDF sum. The recipe is in
\cref{r03:sec:repro-qemu}.
\begin{itemize}
  \item The generated f64 AA kernel with prefetch is bit-identical to the generated f64 pull
        kernel ($N = 16$ and 32). This confirms lbmpy's AA kernels and the driver's ghost-layer
        handling for AA at the same time. f32 differs by about $10^{-8}$, because of FMA
        rounding.
  \item The new instruction selection leaves the f64 checksums unchanged to 13 digits.
  \item With row pointers and \code{in_register}, all three checksums are bit-identical to those
        of the previous kernels, f32 and f64, AA and pull. The kernels without output keep the
        weighted PDF sum of the kernel they derive from: the PDFs are unchanged.
  \item The f64 split kernels agree with the generated AA kernel in all printed digits of all
        three checksums, at $16^3$ cells $\times$ 4 steps and $32^3 \times 10$ steps, for every
        BY, GM and GP tried. f32 differs by $3.5\cdot10^{-7}$, because the sums are formed in a
        different order.
\end{itemize}

\section{Checksums of the board runs}\label{r03:sec:ver-board}

\begin{itemize}
  \item Run~B (\cref{r03:sec:res-runs}), GCC, f64: all seven variants print the checksum
        $3.951592090140\cdot10^{2}$: scalar and RVV, pull and AA, with and without prefetch (log
        \file{e2e-aa-gcc-f64.log}).
  \item Run~B, GCC, f32: pull and AA agree to $8\cdot10^{-8}$ (RVV) and $4\cdot10^{-8}$ (scalar);
        \code{explicit_cse} and \code{explicit_cse_pf} agree exactly, as do
        \code{explicit_cse_aa} and \code{explicit_cse_aa_pf} (log \file{e2e-aa-gcc-f32.log}).
  \item Run~C, f32 and f64: every checksum is identical to run~B's, digit for digit (logs
        \file{e2e-cg-gcc-f32.log}, \file{e2e-cg-gcc-f64.log}).
  \item Acceptance test of the new instruction selection in the lab: checksums of the original
        and the new kernels within $5\cdot10^{-9}$ (f32) and identical (f64) (log
        \file{lab3-gate.log}).
  \item 4 and 8 processes (\cref{r03:sec:res-mpi}): the checksums of all variants differ by at
        most $2.9\cdot10^{-6}$ (f32) and 0 (f64) (log \file{e2e-mpi.log}).
\end{itemize}

\section{Errors found on the way}\label{r03:sec:ver-bugs}

\begin{description}
  \item[Initialisation of the pull benchmark.] The benchmark as it was before this work
        corrupted border PDFs with a halo exchange before the first step
        (\cref{r03:sec:wl-app}). Found because the result depended on the division into blocks.
  \item[A test that could not test.] The first board test of the new instruction selection
        summed all forms into one expression. sympy simplified it, so the \code{vfrsub_vf} that
        the test looked for never appeared, and the test failed before any numbers were
        compared. The guarded session script restored the old instruction selection, as designed
        (\cref{r03:fig:jobflow}). The test now uses one output per form; the second attempt
        passed (145 tests).
  \item[Order of the passes in the split driver.] A first version of \code{lbm_split.py} ran
        the macroscopic pass of the version without output after the collision passes, so the
        collision read the previous block's scratch values. The weighted PDF sum under QEMU
        showed it (7696 instead of 38.72 at $16^3$). Each version now keeps its own order of
        passes.
  \item[lbmpy's equilibrium.] Copying the textbook formula would have generated a slightly
        different model. The symbolic comparison with lbmpy caught it (\cref{r03:sec:wl-split}).
  \item[Measurement artefacts.] The fastest step of an AA kernel is always an even step
        (\cref{r03:sec:lab}). Measured by fastest step, the f32 AA kernel seemed 21\,\% slower in
        the application than in the lab; by cycle average the gap is 15\,\%
        (\cref{r03:tab:labapp}). The errors found in the probes themselves are described at the end of \cref{r03:sec:ops}.
\end{description}
